% Bagian: computability
% Bab: recursive-functions
% Subbagian: pr-relations

\documentclass[../../../include/open-logic-section]{subfiles}

\begin{document}

\olfileid[id]{cmp}{rec}{prr}
\olsection{Relasi Rekursif Primitif}


\begin{defn}
Suatu relasi $R(\vec x)$ disebut rekursif primitif jika fungsi
karakteristiknya,
\[
\Char{R}(\vec x) = \left\{
  \begin{array}{ll}
  1 & \mbox{jika $R(\vec x)$} \\
  0 & \mbox{selain itu}
  \end{array}
\right.
\]
bersifat rekursif primitif.
\end{defn}

Dengan kata lain, ketika kita berbicara tentang relasi rekursif primitif
$R(\vec x)$, yang dimaksud ialah relasi berbentuk $\Char{R}(\vec x) = 1$,
dengan $\Char{R}$ suatu fungsi rekursif primitif yang mengembalikan 1 atau
0 untuk setiap masukan. Sebagai contoh, relasi $\fn{IsZero}(x)$, yang
berlaku jika dan hanya jika $x = 0$, bersesuaian dengan fungsi
$\Char{\fn{IsZero}}$ yang didefinisikan melalui rekursi primitif oleh
\begin{align*}
\Char{\fn{IsZero}}(0) & = 1,\\
\Char{\fn{IsZero}}(x+1) & = 0.
\end{align*}

Jelas bahwa relasi dapat dikomposisikan dengan fungsi-fungsi rekursif primitif
lainnya. Jadi, relasi-relasi berikut juga bersifat rekursif primitif:
\begin{enumerate}
\item Relasi kesamaan, $x = y$, yang didefinisikan oleh
  $\fn{IsZero}(\left|x-y\right|)$.
\item Relasi kurang-dari-atau-sama-dengan, $x \leq y$, yang didefinisikan oleh
  $\fn{IsZero}(x \tsub y)$.
\end{enumerate}


\begin{prop}
  Himpunan relasi rekursif primitif tertutup terhadap operasi Boolean. Artinya,
  jika $P(\vec x)$ dan $Q(\vec x)$ bersifat rekursif primitif, maka demikian
  pula
  \begin{enumerate}
  \item $\lnot P(\vec x)$
  \item $P(\vec x) \land Q(\vec x)$
  \item $P(\vec x) \lor Q(\vec x)$
  \item $P(\vec x) \lif Q(\vec x)$
  \end{enumerate}
\end{prop}

\begin{proof}
  Andaikan $P(\vec x)$ dan $Q(\vec x)$ bersifat rekursif primitif, yakni fungsi
  karakteristiknya $\Char{P}$ dan $\Char{Q}$ bersifat demikian. Kita harus
  menunjukkan bahwa fungsi karakteristik $\lnot P(\vec x)$, dan seterusnya,
  juga bersifat rekursif primitif.
  \[
  \Char{\lnot P}(\vec x) = \begin{cases}
    0 & \text{jika $\Char{P}(\vec x) = 1$}\\
    1 & \text{selain itu}
  \end{cases}
  \]
  Kita dapat mendefinisikan $\Char{\lnot P}(\vec x)$ sebagai
  $1 \tsub \Char{P}(\vec x)$.
  \[
  \Char{P \land Q}(\vec x) = \begin{cases}
    1 & \text{jika $\Char{P}(\vec x) = \Char{Q}(\vec x) = 1$}\\
    0 & \text{selain itu}
  \end{cases}
  \]
  Kita dapat mendefinisikan $\Char{P \land Q}(\vec x)$ sebagai
  $\Char{P}(\vec x) \cdot \Char{Q}(\vec x)$ atau sebagai
  $\fn{min}(\Char{P}(\vec x), \Char{Q}(\vec x))$. Demikian pula,
  \begin{align*}
    \Char{P \lor Q}(\vec x) & = \fn{max}(\Char{P}(\vec x), \Char{Q}(\vec x)) \text{ dan}\\
    \Char{P \lif Q}(\vec x) & = \fn{max}(1 \tsub
  \Char{P}(\vec x), \Char{Q}(\vec x)).
  \end{align*}
\end{proof}

\begin{prop}
  Himpunan relasi rekursif primitif tertutup terhadap kuantifikasi berbatas.
  Artinya, jika $R(\vec x, z)$ merupakan relasi rekursif primitif, maka
  demikian pula relasi
  \begin{align*}
    & \bforall{z < y}{R(\vec x, z)} \text{ dan}\\
    & \bexists{z < y}{R(\vec x, z)}.
  \end{align*}
  $\bforall{z < y}{R(\vec x, z)}$ berlaku bagi $\vec x$ dan $y$ jika dan hanya
  jika $R(\vec x, z)$ berlaku untuk setiap~$z$ yang kurang dari~$y$, dan
  demikian pula untuk $\bexists{z < y}{R(\vec x, z)}$.
\end{prop}

\begin{proof}
  Berdasarkan konvensi, kita menetapkan $\bforall{z < 0}{R(\vec x, z)}$
  bernilai benar (dengan alasan sepele bahwa tidak ada $z$ yang kurang
  dari~$0$), sedangkan $\bexists{z < 0}{R(\vec x, z)}$ bernilai salah.
  Kuantor universal berbatas bekerja persis seperti hasil kali berhingga atau
  minimum berulang. Artinya, jika $P(\vec x, y) \defiff \bforall{z <
  y}{R(\vec x, z)}$, maka $\Char{P}(\vec x, y)$ dapat didefinisikan oleh
  \begin{align*}
    \Char{P}(\vec x, 0) & = 1\\
    \Char{P}(\vec x, y+1) & =
    \fn{min}(\Char{P}(\vec x, y), \Char{R}(\vec x, y)).
  \end{align*}
  Kuantifikasi eksistensial berbatas juga dapat didefinisikan dengan
  $\fn{max}$. Sebagai alternatif, kuantifikasi ini dapat didefinisikan dari
  kuantifikasi universal berbatas dengan menggunakan ekuivalensi
  $\bexists{z < y}{R(\vec x, z)} \liff \lnot \bforall{z < y}{\lnot R(\vec x, z)}$.
  Perhatikan, misalnya, bahwa kuantor berbatas berbentuk
  $\bexists{x \leq y}{\dots x\dots}$ ekuivalen dengan
  $\bexists{x < y+1}{\dots x \dots}$.
\end{proof}

\begin{prob}
  Tunjukkan bahwa relasi berargumen tiga $x \equiv y \mod n$ (kongruensi
  modulo~$n$) bersifat rekursif primitif.
\end{prob}

Fungsi rekursif primitif lain yang berguna ialah fungsi bersyarat
$\fn{cond}(x,y,z)$, yang didefinisikan oleh
\begin{align*}
  \fn{cond}(x,y,z) & = \begin{cases}
  y & \text{jika $x = 0$} \\
  z & \text{selain itu}.
\end{cases}
\intertext{Fungsi ini didefinisikan secara rekursif oleh}
\fn{cond}(0,y,z) & = y,\\
\fn{cond}(x+1,y,z) & = z.
\end{align*}
Kita dapat menggunakannya untuk membenarkan definisi fungsi rekursif primitif
berdasarkan kasus dari relasi rekursif primitif:

\begin{prop}
Jika $g_0(\vec x)$, \dots,~$g_m(\vec x)$ merupakan fungsi rekursif primitif,
dan $R_0(\vec x)$, \dots, $R_{m-1}(\vec x)$ merupakan relasi rekursif
primitif, maka fungsi $f$ yang didefinisikan oleh
\[
f(\vec x) = \begin{cases}
    g_0(\vec x) & \text{jika $R_0(\vec{x})$} \\
    g_1(\vec x) & \text{jika $R_1(\vec{x})$ dan bukan $R_0(\vec{x})$} \\
    \vdots & \\
    g_{m-1}(\vec x) & \text{jika $R_{m-1}(\vec{x})$ dan tidak satu pun kasus
      sebelumnya berlaku}
    \\
    g_m(\vec x) & \mbox{selain itu}
\end{cases}
\]
juga bersifat rekursif primitif.
\end{prop}

\begin{proof}
  Jika $m = 1$, fungsi ini hanyalah fungsi yang didefinisikan oleh
  \[
  f(\vec x) = \fn{cond}(\Char{\lnot R_0}(\vec x),g_0(\vec x),g_1(\vec
  x)).
  \]
  Untuk $m$ yang lebih besar daripada $1$, kita cukup mengomposisikan
  definisi-definisi berbentuk ini.
\end{proof}
\end{document}
